%======================================================================
% MOCKUP SANDRA MAGAZINE - Revista Científica
%======================================================================

\documentclass[12pt, spanish, letterpaper]{book}

\usepackage{fontspec}
\setmainfont{Inter}[Path = assets/fonts/, Extension = .otf, UprightFont = *-Regular]

\usepackage[
    margin=2.5cm,
    top=2.5cm,
    bottom=2.5cm,
    left=3cm,
    right=2cm
]{geometry}

\usepackage{graphicx}
\usepackage{wrapfig}
\usepackage{fancyhdr}
\usepackage{titlesec}
\usepackage{makeidx}
\usepackage{setspace}
\usepackage{multicol}
\usepackage{longtable}
\usepackage{array}
\usepackage{booktabs}
\usepackage{caption}
\usepackage{subcaption}
\usepackage{enumitem}
\usepackage{float}
\usepackage{listings}
\usepackage{xcolor}
\usepackage{lettrine}
\usepackage{epigraph}
\usepackage{microtype}
\usepackage{eso-pic}
\usepackage{tikz}
\usetikzlibrary{calc, shapes.geometric, backgrounds, shadows.blur, arrows.meta}

\usepackage[most]{tcolorbox}
\usepackage{fontawesome5}

%--------------------------------------------------------------------
% COLORES
%--------------------------------------------------------------------
\definecolor{sandraMagDark}{HTML}{1A3A4A}
\definecolor{sandraMagAccent}{HTML}{00A896}
\definecolor{sandraMagGold}{HTML}{D4A574}
\definecolor{sandraMagLight}{HTML}{F4F8F9}
\definecolor{sandraMagGray}{HTML}{E8ECF0}
\definecolor{isoBlue}{HTML}{0066CC}
\definecolor{isoGreen}{HTML}{00875A}
\definecolor{isoOrange}{HTML}{FF8B00}
\definecolor{isoRed}{HTML}{DE350B}

%--------------------------------------------------------------------
% CAJAS
%--------------------------------------------------------------------
\newenvironment{cajadestacada}[1]
{\begin{tcolorbox}[colback=sandraMagLight, colframe=sandraMagAccent, title={\faIcon{lightbulb} #1}, title style={sandraMagDark, font=\bfseries}, boxrule=1pt, rounded corners, after skip=1em]}{\end{tcolorbox}}

\newenvironment{cajaISO}[1]
{\begin{tcolorbox}[colback=sandraMagGray, colframe=isoBlue, title={\faIcon{certificate} #1}, title style={isoBlue, font=\bfseries}, boxrule=2pt, sharp corners, after skip=1em]}{\end{tcolorbox}}

\newenvironment{cajaAdvertencia}
{\begin{tcolorbox}[colback=#FFF3E0, colframe=isoOrange, title={\faIcon{exclamation-triangle} ADVERTENCIA}, title style={isoOrange, font=\bfseries}, boxrule=2pt, rounded corners, after skip=1em]}{\end{tcolorbox}}

\newenvironment{cajaInfo}
{\begin{tcolorbox}[colback=#E3F2FD, colframe=isoBlue, title={\faIcon{info-circle} INFORMACIÓN}, title style={isoBlue, font=\bfseries}, boxrule=1pt, rounded corners, after skip=1em]}{\end{tcolorbox}}

\newenvironment{definicionMagazine}
{\begin{tcolorbox}[colback=sandraMagLight, colframe=sandraMagAccent, title={\faIcon{book} \textbf{Definición}}, title style={sandraMagDark}, boxrule=1pt, rounded corners, after skip=1em, fontupper=\itshape]}{\end{tcolorbox}}

\newenvironment{listaTecnologias}
{\begin{itemize}[itemsep=0.5em, parsep=0.3em, label={\faIcon{cog}}]}{\end{itemize}}

%--------------------------------------------------------------------
% BANNER CAPÍTULO
%--------------------------------------------------------------------
\newcommand{\bannerCapitulo}[3]{%
\begin{tikzpicture}[remember picture, overlay]
    \fill[sandraMagDark] (current page.north west) -- (current page.north east) -- (current page.south east) -- (current page.south west) -- cycle;
    \fill[sandraMagAccent] (current page.north west) rectangle ([yshift=-5mm] current page.north east);
    \node[anchor=west, text=white, font=\fontsize{32}{38}\selectfont\bfseries] at ([xshift=2cm, yshift=-2.5cm] current page.north west) {#1};
    \node[anchor=west, text=sandraMagAccent, font=\fontsize{14}{18}\selectfont] at ([xshift=2cm, yshift=-4cm] current page.north west) {#2};
    \node[anchor=east, text=sandraMagGold, font=\fontsize{72}{86}\selectfont\bfseries] at ([xshift=-2cm, yshift=-3cm] current page.north east) {#3};
    \draw[sandraMagGold, line width=3pt] ([xshift=2cm, yshift=-4.5cm] current page.north west) -- ([xshift=-8cm, yshift=-4.5cm] current page.north east);
\end{tikzpicture}
\clearpage
}

%--------------------------------------------------------------------
% DROP CAP
%--------------------------------------------------------------------
\newcommand{\dropcap}[2]{\lettrine[lines=3, lhang=0.2, findent=0.5em, nindent=0em]{#1}{#2}}

%--------------------------------------------------------------------
% ESTILOS TÍTULOS
%--------------------------------------------------------------------
\titleformat{\chapter}{\normalfont\Huge\bfseries\sandraMagDark}{\thechapter}{1em}{}
\titleformat{\section}{\normalfont\LARGE\bfseries\sandraMagDark}{\thesection}{1em}{\raggedright}

%--------------------------------------------------------------------
% PÁGINAS
%--------------------------------------------------------------------
\pagestyle{fancy}
\fancyhf{}
\fancyhead[R]{\textcolor{sandraMagDark}{\rightmark}}
\fancyhead[L]{\textcolor{sandraMagDark}{\leftmark}}
\fancyfoot[C]{\thepage}

\geometry{headsep=1cm, footskip=1.5cm}

\usepackage[spanish]{babel}

\begin{document}

%======================================================================
% CAPÍTULO MOCKUP - ESTILO REVISTA
%======================================================================
\chapter{Marco Te\'orico}

\bannerCapitulo{Marco Te\'orico}{Fundamentos tecnol\'ogicos y arquitectura de sistemas empresariales}{02}

\section{Estado del Arte}

\dropcap{E}{l} middleware empresarial ha evolucionado significativamente desde sus or\'igenes hasta convertirse en plataformas complejas de integraci\'on. El mercado ofrece soluciones como IBM Integration Bus, MuleSoft, Apache Camel y Spring Integration.

Sin embargo, la mayor\'ia fueron dise\~nadas para mercados desarrollados, ignorando las necesidades espec\'ificas de regiones con infraestructura limitada y contextos regulatorios distintos. Esta brecha crea una oportunidad para soluciones localmente desarrolladas.

\begin{cajaISO}{ISO 27001:2022}
    La familia de normas ISO/IEC 27000 proporciona un marco sistem\'atico para gestionar la seguridad de la informaci\'on. Sandra Ecosystem implementa estos principios en todas las capas.
\end{cajaISO}

\begin{definicionMagazine}
\textbf{Estado del Arte}: Conjunto de tecnolog\'ias, m\'etodos y pr\'acticas m\'as avanzadas para resolver un problema espec\'ifico en un momento determinado.
\end{definicionMagazine}

\section{Fundamentos de Middleware}

El middleware es un software que act\'ua como puente entre aplicaciones y el sistema operativo, facilitando la comunicaci\'on e integraci\'on. Un Enterprise Service Bus (ESB) representa la evoluci\'on hacia una arquitectura basada en servicios.

\begin{cajadestacada}{Implementaci\'on en Sandra Server}
    Sandra Server implementa estos conceptos utilizando Go, proporcionando un ESB de alto rendimiento capaz de manejar miles de conexiones simult\'aneas.
\end{cajadestacada}

\begin{figure}[H]
    \centering
    \begin{tikzpicture}[
        node distance=1.8cm,
        app/.style={rectangle, draw, rounded corners, minimum width=2cm, minimum height=0.8cm, fill=sandraMagLight},
        mid/.style={rectangle, draw, rounded corners, minimum width=3cm, minimum height=1cm, fill=sandraMagAccent!20, font=\bfseries},
        os/.style={rectangle, draw, rounded corners, minimum width=2cm, minimum height=0.8cm, fill=sandraMagGray}
    ]
        \node[app] (app1) {App 1};
        \node[app] (app2) {App 2};
        \node[app] (app3) {App N};
        
        \node[mid, below=1.5cm of app2] (esb) {Enterprise Service Bus};
        
        \node[os, below=1.5cm of esb] (so) {Sistema Operativo};
        
        \draw[->, thick] (app1) -- (esb);
        \draw[->, thick] (app2) -- (esb);
        \draw[->, thick] (app3) -- (esb);
        \draw[->, thick] (esb) -- (so);
    \end{tikzpicture}
    \caption{Arquitectura b\'asica de Middleware}
\end{figure}

\section{Arquitectura de Microservicios}

\dropcap{L}{a} arquitectura de microservicios construye aplicaciones como servicios peque\~nos, aut\'onomos y dbilmente acoplados. Cada microservicio es responsable de una funcionalidad espec\'ifica.

\begin{cajaInfo}
El ecosistema Sandra implementa esta arquitectura:
\end{cajaInfo}

\begin{listaTecnologias}
    \item \textbf{Server}: Middleware central de comunicaciones
    \item \textbf{Consola}: Interfaz administrativa web
    \item \textbf{SDC}: Aplicaci\'on de escritorio
    \item \textbf{Sentinel}: Procesamiento especializado
\end{listaTecnologias}

\begin{cajaAdvertencia}
    La migraci\'on a microservicios requiere planificaci\'on cuidadosa. Se recomienda comenzar con servicios independientes.
\end{cajaAdvertencia}

\begin{wrapfigure}{r}{0.4\textwidth}
    \centering
    \begin{tikzpicture}[
        node distance=0.8cm,
        ms/.style={circle, draw, fill=sandraMagAccent!30, minimum width=1.2cm, font=\bfseries\tiny}
    ]
        \node[ms] (s1) {Server};
        \node[ms] (s2) [right of=s1] {Consola};
        \node[ms] (s3) [below of=s1] {SDC};
        \node[ms] (s4) [below of=s2] {Sentinel};
        
        \draw[<->, thick] (s1) -- (s2);
        \draw[<->, thick] (s1) -- (s3);
        \draw[<->, thick] (s1) -- (s4);
    \end{tikzpicture}
    \caption{Ecosistema Sandra}
\end{wrapfigure}

\section{Seguridad Zero Trust}

\dropcap{E}{l} modelo Zero Trust elimina la confianza impl\'icita, verificando cada solicitud independientemente de su origen:

\begin{cajaISO}{ISO 27002:2022}
    Los controles de seguridad incluyen 93 medidas organizadas en cuatro categor\'ias: organizacionales, personales, f\'isicos y tecnol\'ogicos.
\end{cajaISO}

\begin{enumerate}
    \item \textbf{Verificaci\'on expl\'icita}: Autenticaci\'on multifactor
    \item \textbf{M\'inimo privilegio}: Permisos m\'inimos necesarios
    \item \textbf{Verificaci\'on continua}: Monitoreo de sesiones
\end{enumerate}

\section{Tecnolog\'ias Emergentes}

\begin{tabular}{m{3cm} m{10cm}}
    \textcolor{sandraMagAccent}{\faIcon{code}} \textbf{Go 1.24} & Alto rendimiento y concurrencia \\[0.6em]
    \textcolor{sandraMagGold}{\faIcon{shield-alt}} \textbf{Rust 2021} & Seguridad de memoria \\[0.6em]
    \textcolor{isoRed}{\faIcon{layer-group}} \textbf{Angular 16+} & Framework frontend \\[0.6em]
    \textcolor{sandraMagDark}{\faIcon{desktop}} \textbf{Tauri 2.0} & Apps de escritorio \\[0.6em]
    \textcolor{isoGreen}{\faIcon{network-wired}} \textbf{gRPC} & Comunicaci\'on entre servicios \\[0.6em]
    \textcolor{isoBlue}{\faIcon{bolt}} \textbf{WebSockets} & Tiempo real bidireccional
\end{tabular}

\begin{cajadestacada}{Nota de Rendimiento}
    Go + Tauri proporciona equilibrio \'optimo entre rendimiento y consumo de recursos.
\end{cajadestacada}

\end{document}
